\chapter{Axiom $\I_1$: Existence of the Indeterminate}
\label{ch:axiom-I1}

\section{Introduction}

The first and most fundamental axiom of our extended system posits the existence of a unique element $\I$ that arises from self-subtraction of cardinals. This axiom is the cornerstone upon which the entire Algebra of the Indeterminate rests.

\section{Motivating the Axiom}

\subsection{The Problem of $\kappa - \kappa$}

In classical cardinal arithmetic, the subtraction $\kappa - \kappa$ for an infinite cardinal $\kappa$ is undefined or, by convention, taken to be zero. However, this treatment obscures important mathematical phenomena:

\begin{enumerate}
    \item \textbf{Limit indeterminacy}: Consider the limit $\lim_{n \to \infty} (a_n - a_n)$ where $(a_n)$ is unbounded. Different interpretations yield different results.
    
    \item \textbf{Ultrapower ambiguity}: In ultrapower constructions, sequences $(f_i)$ where $f_i = g_i$ for almost all $i$ still may not produce determinate differences.
    
    \item \textbf{Algebraic incompleteness}: Setting $\kappa - \kappa = 0$ forces an artificial closure that prevents natural extensions of arithmetic.
\end{enumerate}

\subsection{Historical Precedents}

The notion of indeterminate forms has deep roots in analysis:

\begin{itemize}
    \item \textbf{L'Hôpital's Rule} (1696): Addressed $\frac{0}{0}$ through limits
    \item \textbf{Cauchy's infinitesimals} (1821): Introduced quantities between zero and finite values
    \item \textbf{Robinson's hyperreals} (1961): Formalized infinitesimals rigorously
    \item \textbf{Wheel theory} (1997): Carlström's algebraic structure allowing $0/0$
\end{itemize}

Our approach synthesizes these ideas into a set-theoretic framework.

\section{Formal Statement of Axiom $\I_1$}

\begin{axiom}[Existence of $\I$ -- Axiom $\I_1$]
\label{ax:I1}
There exists a unique element $\I$ such that for every cardinal $\kappa$:
\[
\kappa \ominus \kappa = \I
\]
where $\ominus$ denotes extended cardinal subtraction.
\end{axiom}

\subsection{Unpacking the Formalization}

The axiom contains several components that require clarification:

\begin{definition}[Extended Cardinal Subtraction]
The operation $\ominus : \text{Card} \times \text{Card} \to \text{Card} \cup \{\I\}$ extends standard cardinal subtraction by:
\[
\kappa \ominus \lambda = \begin{cases}
    \mu & \text{if } \kappa = \lambda + \mu \text{ in standard arithmetic} \\
    \I & \text{if } \kappa = \lambda \\
    \text{undefined} & \text{if } \kappa < \lambda
\end{cases}
\]
\end{definition}

\begin{remark}
The case $\kappa < \lambda$ may later be extended to introduce negative cardinals or additional indeterminate elements, but we defer this to future work.
\end{remark}

\section{Proof of Independence}

\subsection{Independence from ZF + AC}

\begin{theorem}[Independence of Axiom $\I_1$]
\label{thm:I1-independence}
Axiom $\I_1$ is independent of ZF + AC. That is, there exist models of ZF + AC where $\I_1$ holds and models where $\I_1$ fails.
\end{theorem}

\begin{proof}
\textbf{Model where $\I_1$ holds:}

We construct a model $\mathcal{M}^+$ by extending a standard model $\mathcal{M}$ of ZF + AC:

\textbf{Step 1}: Start with $\mathcal{M} \models \text{ZF} + \text{AC}$.

\textbf{Step 2}: Introduce a new urelements $\I$ not in the original domain. Define the extended universe:
\[
M^+ = M \cup \{\I\}
\]

\textbf{Step 3}: Extend the operation $\ominus$ by fiat:
\[
\kappa \ominus \lambda = \begin{cases}
    \kappa - \lambda & \text{if } \kappa > \lambda \text{ (standard)} \\
    \I & \text{if } \kappa = \lambda \\
    \text{error} & \text{otherwise}
\end{cases}
\]

\textbf{Step 4}: Verify that all ZF + AC axioms still hold in $\mathcal{M}^+$. Since $\I$ is added as an atom and ZF + AC does not quantify over operations, the original axioms remain satisfied.

\textbf{Step 5}: By construction, $\mathcal{M}^+ \models \I_1$.

\vspace{0.5cm}

\textbf{Model where $\I_1$ fails:}

Consider the standard model $\mathcal{M}$ where cardinal subtraction is defined only for $\kappa > \lambda$ and $\kappa \ominus \kappa$ is simply taken to be $0$ (or undefined). This model satisfies ZF + AC but violates $\I_1$ since there is no element $\I$ satisfying the required property.

Therefore, $\I_1$ is independent of ZF + AC.
\end{proof}

\subsection{Forcing Construction}

We can also establish independence using forcing:

\begin{theorem}[Forcing Extension]
\label{thm:I1-forcing}
Given a model $\mathcal{M}$ of ZF + AC, there exists a forcing extension $\mathcal{M}[G]$ where $\I_1$ holds.
\end{theorem}

\begin{proof}[Proof Sketch]
Define a partial order $\mathbb{P}$ consisting of finite approximations to the operation $\ominus$ that eventually set $\kappa \ominus \kappa = \I$ for a generic element $\I$.

Let $G$ be a $\mathbb{P}$-generic filter over $\mathcal{M}$. In the extension $\mathcal{M}[G]$:
\begin{itemize}
    \item Define $\I = \bigcup \{p(\kappa, \kappa) : p \in G\}$
    \item Extend $\ominus$ to all cardinals using the generic filter
    \item Verify that $\I_1$ holds by the genericity of $G$
\end{itemize}

The standard arguments for forcing ensure that ZF + AC axioms are preserved.
\end{proof}

\section{Consistency Results}

\subsection{Relative Consistency}

\begin{theorem}[Relative Consistency of $\I_1$]
\label{thm:I1-consistency}
If ZF + AC is consistent, then ZF + AC + $\I_1$ is consistent.
\end{theorem}

\begin{proof}
The proof follows from Theorem~\ref{thm:I1-independence}. Since we can construct a model of ZF + AC + $\I_1$ from any model of ZF + AC, a contradiction in the extended theory would project to a contradiction in the base theory.

Formally, suppose $\text{ZF} + \text{AC} + \I_1 \vdash \bot$. Then by the soundness of first-order logic, there is no model of ZF + AC + $\I_1$. But we constructed such a model $\mathcal{M}^+$, a contradiction.
\end{proof}

\subsection{Equiconsistency}

\begin{corollary}[Equiconsistency]
The theories ZF + AC and ZF + AC + $\I_1$ are equiconsistent.
\end{corollary}

\begin{proof}
By Theorem~\ref{thm:I1-consistency}, if ZF + AC is consistent, so is ZF + AC + $\I_1$. Conversely, ZF + AC is a subtheory, so if ZF + AC + $\I_1$ is consistent, then so is ZF + AC.
\end{proof}

\section{Counterexamples: What If $\I_1$ Is Removed?}

\subsection{Loss of Closure}

Without Axiom $\I_1$, the extended cardinal operations lack closure:

\begin{example}[Non-Closed Arithmetic]
Consider the sequence of operations:
\begin{align*}
    (\aleph_0 + \aleph_0) \ominus \aleph_0 &= \aleph_0 \\
    (\aleph_0 \cdot 2) \ominus \aleph_0 &= \aleph_0 \\
    \aleph_0 \ominus \aleph_0 &= \, ?
\end{align*}
Without $\I_1$, the last operation is undefined, breaking the algebraic structure.
\end{example}

\subsection{Failure of Universal Properties}

\begin{proposition}[Loss of Universality]
Without $\I_1$, there is no element $x$ satisfying $x = \kappa \ominus \kappa$ for all cardinals $\kappa$ simultaneously.
\end{proposition}

\begin{proof}
If we try to define separate elements $x_\kappa = \kappa \ominus \kappa$ for each $\kappa$, we have infinitely many distinct elements (one per cardinal), which violates the parsimony principle and creates an unnecessarily complex structure.

Alternatively, if we set $x_\kappa = 0$ for all $\kappa$, we obtain a trivial theory (as shown in Chapter~\ref{ch:preliminary-logic}, Theorem~\ref{thm:triviality-collapse}).
\end{proof}

\section{Properties Derived from Axiom $\I_1$}

\subsection{Immediate Consequences}

\begin{proposition}[Self-Indeterminacy]
For any cardinal $\kappa$:
\[
(\kappa \ominus \kappa) \ominus (\kappa \ominus \kappa) = \I \ominus \I = \I
\]
\end{proposition}

\begin{proof}
By $\I_1$, $\kappa \ominus \kappa = \I$. Applying $\I_1$ again with $\lambda = \I$ (treating $\I$ as an element of the extended domain):
\[
\I \ominus \I = \I
\]
This will be formalized as Axiom $\I_5$ in Chapter~\ref{ch:axiom-I5}.
\end{proof}

\begin{proposition}[Idempotence Class]
The element $\I$ belongs to the class of idempotent elements under $\ominus$.
\end{proposition}

\subsection{Interaction with Cardinal Arithmetic}

\begin{theorem}[Cardinal Sum Behavior]
For cardinals $\kappa, \lambda$:
\[
(\kappa + \lambda) \ominus (\kappa + \lambda) = \I
\]
\end{theorem}

\begin{proof}
Let $\mu = \kappa + \lambda$. Since $\mu$ is a cardinal, by Axiom $\I_1$:
\[
\mu \ominus \mu = \I
\]
Substituting back:
\[
(\kappa + \lambda) \ominus (\kappa + \lambda) = \I
\]
\end{proof}

\section{Connection to Other Axioms}

Axiom $\I_1$ serves as the foundation for the remaining axioms:

\begin{itemize}
    \item \textbf{Axiom $\I_2$}: Extends $\I_1$ to cases where $\kappa \neq \lambda$
    \item \textbf{Axiom $\I_4$}: Uses $\I$ to define absorption properties
    \item \textbf{Axiom $\I_5$}: Formalizes $\I \ominus \I = \I$ hinted at above
    \item \textbf{Axiom $\I_7$}: Ensures $\I$ preserves non-triviality
\end{itemize}

The interplay between these axioms will be explored in subsequent chapters.

\section{Philosophical Considerations}

\subsection{Ontological Status of $\I$}

Axiom $\I_1$ raises the question: \textit{What is $\I$?}

\begin{itemize}
    \item \textbf{Platonist view}: $\I$ is an abstract object discovered through mathematical investigation, existing independently of our definitions.
    
    \item \textbf{Formalist view}: $\I$ is a symbol introduced axiomatically, with meaning derived solely from its role in the formal system.
    
    \item \textbf{Intuitionist view}: $\I$ represents the failure of a mental construction, a state where our computational process does not terminate with a determinate value.
\end{itemize}

We adopt a pragmatic stance: $\I$ is as real as any mathematical object, justified by its utility and consistency.

\subsection{Comparison to Other Indeterminates}

How does $\I$ compare to other "undefined" or "indeterminate" concepts?

\begin{center}
\begin{tabular}{|l|l|l|}
\hline
\textbf{Concept} & \textbf{Context} & \textbf{Relation to $\I$} \\
\hline
$0/0$ & Real analysis & $\I$ formalizes this \\
$\infty - \infty$ & Extended reals & Special case: $\aleph_0 \ominus \aleph_0 = \I$ \\
NaN & Computation & $\I$ replaces NaN with algebraic structure \\
$\bot$ & Logic & $\I$ is truth value, not falsehood \\
Infinitesimals & NSA & $\I$ generalizes to cardinals \\
\hline
\end{tabular}
\end{center}

\section{Computational Interpretation}

\subsection{Implementation in Python}

The following code snippet demonstrates how $\I$ can be implemented:

\begin{verbatim}
class IndeterminateElement:
    """Represents the indeterminate element I."""
    
    def __init__(self):
        self.value = "I"
    
    def __sub__(self, other):
        if isinstance(other, IndeterminateElement):
            return IndeterminateElement()  # I - I = I
        # Handle other cases...
    
    def __repr__(self):
        return "I"

# Axiom I1 in action
def cardinal_self_subtract(kappa):
    """For any cardinal κ, returns κ ⊖ κ = I."""
    return IndeterminateElement()
\end{verbatim}

\subsection{Type-Theoretic Formulation}

In dependent type theory (e.g., Lean, Agda):

\begin{verbatim}
inductive ExtendedCardinal : Type
  | card : Cardinal → ExtendedCardinal
  | indet : ExtendedCardinal

def ominus : ExtendedCardinal → ExtendedCardinal → ExtendedCardinal
  | (card κ), (card κ) => indet
  | (card κ), (card λ) => if κ > λ then card (κ - λ) else error
  | indet, indet => indet
  | _, _ => error

axiom I1 : ∀ κ : Cardinal, ominus (card κ) (card κ) = indet
\end{verbatim}

\section{Open Questions}

\begin{enumerate}
    \item Can $\I_1$ be weakened to a weaker existence axiom without losing essential properties?
    
    \item Is there a natural categorical interpretation of $\I$ as a terminal object in a category of "partial cardinals"?
    
    \item How does $\I_1$ interact with large cardinal axioms (inaccessibles, measurables, etc.)?
    
    \item Can we extend $\I_1$ to ordinal arithmetic, defining $\alpha \ominus \alpha = \I$ for ordinals $\alpha$?
\end{enumerate}

\section{Summary}

This chapter has established Axiom $\I_1$ as the foundational postulate of our extended system:

\begin{itemize}
    \item \textbf{Statement}: $\kappa \ominus \kappa = \I$ for all cardinals $\kappa$
    \item \textbf{Independence}: $\I_1$ is independent of ZF + AC
    \item \textbf{Consistency}: ZF + AC + $\I_1$ is equiconsistent with ZF + AC
    \item \textbf{Consequences}: Establishes $\I$ as a universal indeterminate element
    \item \textbf{Necessity}: Without $\I_1$, extended cardinal arithmetic lacks closure
\end{itemize}

With this foundation in place, we proceed to Axiom $\I_2$, which governs the behavior of $\ominus$ for distinct cardinals.

\section{Exercises}

\begin{enumerate}
    \item Verify that in the model $\mathcal{M}^+$, all axioms of ZF remain satisfied after adjoining $\I$.
    
    \item Show that if we allow multiple indeterminate elements $\I_1, \I_2, \ldots$, we lose uniqueness and the theory becomes unnecessarily complex.
    
    \item Prove that Axiom $\I_1$ cannot be derived from any finite subset of ZF + AC axioms.
    
    \item Construct an explicit forcing poset $\mathbb{P}$ for Theorem~\ref{thm:I1-forcing}.
    
    \item Explore the consequences of defining $\kappa \ominus \kappa = \kappa$ instead of $\I$. Show that this leads to contradictions or triviality.
\end{enumerate}
